\scanpage{5}
\mspage{4}
strates de $F$. \deleted{Donc on peut \unsure{ide|idé}}
\[
\begin{aligned}
 F &\quad \text{figure}\\
 \Sigma_F &\quad \text{ens de ses strates}\\
 \Sigma'_F &\quad \text{ens des sous-figures élémentaires.}
\end{aligned}
\]
On :
\diagram{\begin{tikzcd}[ampersand replacement=\&,column sep=small,row sep=.35em]
\Sigma'_F \& {}\simeq\Sigma_F\\
{\unsure{\rotatebox{-90}{$\subset$}|\cap|\rotatebox{-90}{$\in$}}} \&\\
\mathcal{F}_{\text{él}} \&
\end{tikzcd}}
$F$ est connu quand on connait $\Sigma'_F\subset\mathcal{F}_{\text{él}}$ car $F$ est l'ens des strates des $F'\in\Sigma'_F$.

Mais la relation d'inclusion entre figures \unsure{s'interprete|s'interprète} comme relation d'inclusion entre \unsure{partie|parties} de $\mathcal{F}_{\text{él}}$.

\deleted{Deux f}

Considérons la relation d'inclusion \inserted{$\leq$} entre figures élémentaires. \deleted{\unsure{A}}Alors, $\mathcal{F}_{\text{él}}$ est un ens ordonné. \deleted{Deux} \inserted{Si} \deleted{figures sont dites} \inserted{$F$ est une figure,} $\Sigma'_F$ est une partie fermée de l'ens ordonné $\mathcal{F}_{\text{él}}$. L'\unsure{inclussion|inclusion} entre figures est l'inclusion entre ces parties fermées.
